\documentclass[11pt]{article}
\usepackage[a4paper,margin=2.3cm]{geometry}
\usepackage[T1]{fontenc}
\usepackage{mathptmx}
\usepackage{amsmath,amssymb}
\usepackage{xcolor}
\usepackage{enumitem}
\usepackage{xr}
\externaldocument{main}
\usepackage[hidelinks]{hyperref}
\setlength{\parindent}{0pt}
\setlength{\parskip}{5pt}
\newcommand{\rc}[2]{\par\medskip\noindent\colorbox{black!6}{\parbox{\dimexpr\linewidth-2\fboxsep}{\textbf{Comment #1.} \itshape #2}}\par}
\newcommand{\resp}{\par\noindent\textbf{Response.} }
\newcommand{\chg}{\par\noindent\textbf{Changes.} }
\hypersetup{pdftitle={Response to reviewers, algorithms-4716513},pdfauthor={Md. Redwan Islam}}

\begin{document}
\begin{center}
{\large\textbf{Response to the Reviewers}}\\[3pt]
Manuscript algorithms-4716513, ``Event-Aware AGN-Net: Learning Which Graph Updates Can Be Skipped in Streaming Respiratory Sound Event Detection''
\end{center}

Dear Editor, dear Reviewers,

Thank you for the careful and constructive reviews. I have revised the manuscript accordingly and carried out the experiments the reviewers asked for. Changes are highlighted in the accompanying tracked-changes PDF; the comments are numbered R1.1--R1.9 and R2.1--R2.3 below. In summary, the revision adds:
\begin{itemize}[leftmargin=1.3em,itemsep=1pt]
\item a gate trained on the task loss with straight-through Gumbel selection and the same budget-exact rule, as a new baseline, on two cells with three seeds each (R1.3, R2.1);
\item a second dataset, the adult HF\_Lung\_V1 database, with the full pipeline for two cells (R1.4);
\item a harmonic graph topology that links each frequency region to its second and third harmonics (R2.2);
\item an analysis of one-step versus multi-step skip harm (R1.6), long patient streams with a budget-adaptive refresh rule (R1.7), measured contraction ratios and a state-to-output link for Proposition~1 (R1.8), boundary-tolerant metrics (R1.1), and a scaling analysis of the controller overhead (R1.2);
\item a re-framed narrative around the model-agnostic selection method (R1.5), a stage-by-stage statement of the use of generative AI, and the public release of the code, logs, per-recording results and weights (R2.3).
\end{itemize}
Several new results are mixed or negative, and we report them as such. All model results quoted below, apart from a few rounded summary comparisons, are copied automatically from the result files, as in the manuscript.

\section*{Reviewer 1}

\rc{R1.1}{Due to the annotation style shift between the training and test releases, absolute event-based F-scores remain very low across all models. Emphasize more prominently in the abstract and early sections that the primary evaluation target is model preservation ($F_{\mathrm{pres}}$) rather than absolute benchmark accuracy, and consider evaluating with boundary-tolerant metrics.}
\resp We agree. The abstract, the introduction and the opening of the Experiments section now state that the primary target is preservation of the full model's own detections and that accuracy against the annotations is secondary because of the annotation shift. We added boundary-tolerant metrics (Table~\ref{tab:tolerant}): event $F$ with 0.5\,s and 1\,s collars, onset-only and overlap matching, and segment-based $F$ with 1\,s and 2\,s segments, against the annotations and for preservation. For the full-update model on the 2024 test set, event $F$ rises from 0.092 with the strict 0.2\,s collars to 0.308 with 1\,s collars and 0.633 with overlap matching, so much of the low strict-collar score comes from boundary conventions. In preservation, our controller is ahead of the end-to-end gate, the change objective and random selection under every criterion; against the annotations it is ahead of the change objective and random selection under every event-based criterion. We also report one result that cuts the other way: against the annotations, the end-to-end gate scores slightly higher than our controller under the event-based criteria at 25\% on the 2024 set (lower in adventitious segment $F$), although it preserves far fewer of the full model's events (Section~\ref{sec:results}).
\chg Abstract; Section~\ref{sec:intro}; Section~\ref{sec:experiments} (opening) and Section~\ref{sec:metrics}; Section~\ref{sec:results} (``Accuracy against the annotations'', ``Boundary-tolerant metrics'', Table~\ref{tab:tolerant}).

\rc{R1.2}{At a 25\% update budget, the end-to-end speedup is only 1.14$\times$ in NumPy (and slower than full updating in PyTorch at 0.88$\times$) because the controller consumes roughly half of the computation time saved by skipping graph nodes. Include a brief quantitative scaling analysis or discussion on how this overhead diminishes for larger graph cells or dedicated embedded hardware accelerators (DSP/NPU).}
\resp We added a scaling analysis (Table~\ref{tab:scaling}). In the NumPy runtime, the controller's share of the time saved by skipping 75\% of the graph update falls from 60.4\% for our cell ($d{=}96$) to 19.1\% at $d{=}384$ and 8.9\% at $d{=}768$ (7.1\% with four layers), and the graph-stage speedup rises from 1.25$\times$ to 1.82$\times$. The speedup levels off well below the ideal $4\times$, because in this runtime updating a quarter of the nodes still takes about half the time of a full graph update at every size; we state this explicitly. For embedded accelerators we give a qualitative discussion (the controller is a small dense network that suits a DSP or NPU, while sparse gather/scatter of a changing node set does not) and state that we did not measure on such hardware.
\chg Section~\ref{sec:analyses} (``Overhead scaling''); Section~\ref{sec:h3} (``How the overhead scales'', Table~\ref{tab:scaling}).

\rc{R1.3}{The study omits end-to-end differentiable gating baselines (such as Skip-RNN or Gumbel-Softmax token/node pruning). Add at least one end-to-end budgeted gating baseline or provide a deeper empirical justification for why two-stage counterfactual regression outperforms joint optimization.}
\resp We added a gate trained on the task loss of stage C (cross-entropy plus distillation), with the same architecture, inputs and size as our controller and the same budget-exact greedy rule in training and test, using straight-through gradients through Gumbel-perturbed selections. We state its limitations openly: the backbone is frozen while the gate is trained (unlike Skip RNN), stage C fine-tuned the backbone with the utility heads but not with this gate, and its learning rate and temperatures were not tuned, so the comparison may favour our controller to an unknown degree. In preservation, our controller is better in every seed and at every budget: at 25\%, $F_{\mathrm{pres}}=0.775$ against 0.595 for the gate with the AGN-Net cell, and 0.712 against 0.358 with the graph GRU. Against the annotations, with the AGN-Net cell the gate is slightly higher at 25\% on the 2024 set and at 12.5\% on the 2025 release and lower at 50\%; with the graph GRU it is lower at 25\% and 50\%. As empirical justification we report that the gate's task loss does not decrease during training, that its scores rank regions by measured skip-harm only weakly (Spearman 0.24 versus 0.63 for our controller), and that the counterfactual target supervises every region at every step, whereas the gate receives one scalar loss per crop through a biased straight-through estimate of a discrete, budget-exact selection.
\chg Section~\ref{sec:related-cc}; Section~\ref{sec:policies} (``End-to-end gate''); Section~\ref{sec:results} (``Comparison with end-to-end training'', ``Accuracy against the annotations''); Section~\ref{sec:why-two-stage}; Tables~\ref{tab:policy}, \ref{tab:policygg} and~\ref{tab:tolerant}.

\rc{R1.4}{All experiments are conducted exclusively on the SPRSound paediatric dataset. Validate the approach on a second respiratory dataset (e.g., HF\_Lung or ICBHI adapted for streaming) or explicitly narrow the scope claims to paediatric auscultation.}
\resp We added HF\_Lung\_V1 (adult patients, two recording devices, inhalation/exhalation and continuous and discontinuous adventitious sound labels) with the full pipeline for the AGN-Net cell and the graph GRU (three seeds each). Recordings are not labelled by patient, so the validation split is disjoint in pseudo-subjects (recording device and date, as the dataset documentation recommends). The findings carry over: at 25\%, the skip-harm controller preserves $F_{\mathrm{pres}}=0.736$ against 0.638 for the change objective and 0.330 for random selection with the AGN-Net cell, and 0.744 against 0.494 and 0.314 with the graph GRU, better in every seed and at every budget; utility is again concentrated (92.2\% in the top quarter of regions). Against the annotations, non-inferiority to full updating is narrowly missed at 25\% for the AGN-Net cell, missed for the graph GRU, and holds at 50\% for both. During this experiment the AGN-Net cell diverged in three training attempts (the graph GRU never did); the manuscript describes this and the divergence rule we added, which did not trigger in the reported runs, and the logs of all attempts are released.
\chg Section~\ref{sec:data} (HF\_Lung, Table~\ref{tab:data}); Section~\ref{sec:impl} (schedules and divergence rule); Section~\ref{sec:generality} (``Another dataset'', Table~\ref{tab:hflung}); Abstract and Section~\ref{sec:conclusion}.

\rc{R1.5}{The custom acoustic AGN-Net cell underperforms a plain graph-GRU backbone and exhibits high variance across random seeds. Re-frame the architectural narrative around the model-agnostic nature of the skip-harm controller rather than pitching AGN-Net as a superior acoustic backbone.}
\resp We agree and re-framed the paper. The abstract, the contributions and the Method section now present the skip-harm target, the controller and the budget-exact selection rule as a model-agnostic selection method, demonstrated with two cells (the AGN-Net-style cell and a plain graph GRU, whose selective update is now also defined in the Method section), two graph topologies and two datasets. The AGN-Net cell is presented only as one instance that carries over AGN-Net's localized updates, and we keep the explicit statements that it is not a better detector than the graph GRU and that it is less stable to train.
\chg Abstract; Section~\ref{sec:intro}; Section~\ref{sec:method}; Sections~\ref{sec:generality} and~\ref{sec:backbones}; Section~\ref{sec:conclusion}.

\rc{R1.6}{The skip-harm utility is computed using one-step counterfactuals, which ignores the compounding multi-step error propagation caused by leaving a node state stale over consecutive chunks. Discuss this myopic limitation more thoroughly in Section~3.5 and explore multi-step rollout utilities or provide empirical evidence showing why one-step proxies suffice.}
\resp We now discuss the limitation in Section~\ref{sec:utility} and measured multi-step harm directly (Table~\ref{tab:multistep}): each region is held for $k\in\{1,5,25\}$ consecutive steps and the output divergence is summed over horizons of up to 25 steps (1\,s). A skip has lasting effects (over 1\,s, the harm of a single skip is 13.6 times its harm in the step of the skip), but the ranking of regions, which is what the selection uses, is largely stable: the one-step utility ranks regions by their 1\,s harm with Spearman 0.72 and by the harm of 25-step holds with 0.64; the trained controller keeps 0.61, whereas the change head stays at 0.16. The graph GRU shows the same pattern. We also state the limit of this test: it holds one region at a time, whereas in deployment many regions are stale together.
\chg Section~\ref{sec:utility}; Section~\ref{sec:analyses} (``Multi-step harm''); Section~\ref{sec:multistep} (Table~\ref{tab:multistep}); Section~\ref{sec:conclusion} (limitations).

\rc{R1.7}{The forced refresh rule for states older than 50 steps (2 s) is arbitrary and actually degrades $F_{\mathrm{pres}}$ at tight budgets (12.5\%) on short recordings without showing clear empirical benefit. Either make the staleness threshold adaptive to the budget or evaluate the model on longer continuous audio streams where state drift actually occurs.}
\resp We did both. We built 60--120\,s streams by concatenating each test patient's recordings (141 streams, 3.6\,h, from the 2024 test set) and propose a budget-adaptive threshold $A(b)=\lceil 12.5/b\rceil$ that caps the refresh share at 8\% of the budget (Table~\ref{tab:long}). The effects are small and seed-dependent, and we say so. At 25\%, refresh helps on average (0.771 without refresh, 0.793 with the 2\,s rule; strongly in one seed and negligibly in a second, while in the third it lowers preservation by about 0.02), and without refresh the frame-level agreement with the full model drifts down later in the stream, while with refresh it stays level. As the reviewer suspected, the fixed 2\,s threshold is too aggressive at 12.5\% (0.573 versus 0.595 without refresh). No fixed threshold is best at every budget; the adaptive rule, like a fixed $A{=}100$, stays within 0.01 of the best rule at every budget, so we recommend a budget-dependent threshold. We also report that on short recordings the rule costs more for the graph GRU, that the long-stream tests use the AGN-Net cell only, and that two-minute streams are still short of continuous wearable monitoring.
\chg Section~\ref{sec:selection}; Section~\ref{sec:analyses} (``Long streams''); Section~\ref{sec:long} (Table~\ref{tab:long}); Section~\ref{sec:conclusion}.

\rc{R1.8}{Proposition~1 relies on a Lipschitz continuity assumption that is neither enforced during training nor empirically verified for the learned network weights, and it bounds state error rather than readout classification error. Empirically measure the spectral norms/Lipschitz constants of the trained cells or tighten the connection between the state-error bound and the KL-based skip-harm objective.}
\resp We did both, and the results are partly negative. (i) The proof only needs the secant ratio $\gamma_t$ between the two trajectories that actually occur, which we measured on real streams (Table~\ref{tab:lipschitz}). At 25\% its median is 0.970 and 99.2\% of steps contract, but on a few steps the ratio exceeds one (maximum 1.82), so the uniform condition of the geometric bound does not hold; we therefore state the proposition's time-varying form as the main result. That form holds by construction when evaluated with the measured $\gamma_t$, and on the trajectories it is loose. The local Jacobian spectral norm in the Euclidean norm (median 1.29) suggests that the contraction is a property of the trajectories rather than a uniform property of the cell. (ii) We added an explicit link to the output, Eq.~(\ref{eq:readout}): if the readout has Lipschitz constant $L_R$ from states to frame logits, the frame KL divergence is at most $2L_Re_t$. Empirically this link is loose (median readout ratio 0.142, 99th percentile 0.489; state error and KL are nearly uncorrelated, Spearman 0.05), which is why we train the controller on the output-level skip-harm utility rather than on a state-level quantity.
\chg Section~\ref{sec:theory} (Proposition~\ref{prop:bound} and Eq.~(\ref{eq:readout})); Section~\ref{sec:analyses}; Section~\ref{sec:lip} (Table~\ref{tab:lipschitz}).

\rc{R1.9}{Figures and tables can be improved.}
\resp All tables use at least 8\,pt type and no vertical rules; multi-panel figures have panel labels that the captions describe; figure text was enlarged for legibility; and the new analyses are summarised in compact tables.

\section*{Reviewer 2}

\rc{R2.1}{Lack of comparison with end-to-end trainable mechanisms (Skip RNN, Adaptive Computation Time). Without a direct comparison, it cannot be claimed that the proposed two-stage approach genuinely outperforms established methods for conditional computation.}
\resp We added this comparison; please see our response to R1.3. In short, a gate with the same architecture, trained on the task loss with exact budgets and straight-through gradients, is a competitive baseline. It stays below the two-stage controller in preservation of the full model's detections in every seed, at every budget and on both cells; against the annotations it is slightly higher at 25\% on the 2024 set with the AGN-Net cell (but lower with the graph GRU at 25\% and 50\%); and on the graph-GRU cell its results vary strongly between seeds. We now describe its limitations (frozen backbone, untuned settings) and have narrowed the related-work statement: we claim an advantage over the gate we implemented, not over every conditional-computation method.
\chg Sections~\ref{sec:related-cc}, \ref{sec:policies}, \ref{sec:results} and~\ref{sec:why-two-stage}.

\rc{R2.2}{Overly simplistic and static graph topology: a graph that links only adjacent frequency bins cannot propagate information between harmonics, e.g., of wheezes.}
\resp We implemented a harmonically informed topology that additionally links each region to the regions nearest to two and three times its centre frequency (82 extra edges) and ran the full pipeline on it (three seeds). The selection method works unchanged on this graph and stays well ahead of the change objective and random selection (Table~\ref{tab:policyh}): at 25\%, 0.554 against 0.327 and 0.212. The harmonic graph also helps the cell itself: with full updates, the AGN-Net cell's event $F$ rises from 0.087 on the chain graph to 0.139, and the large seed-to-seed variance of the chain-graph cell disappears, bringing it to the level of the plain graph GRU (0.133). We are careful about the interpretation: adventitious event $F$ does not improve, and the harmonic graph is also denser, so the gain cannot be attributed specifically to propagation between harmonics. Absolute preservation is lower on the harmonic graph, and the controller is further from the hindsight oracle there. An adaptive or learned topology remains an interesting extension, to which the controller applies unchanged. We thank the reviewer for this suggestion.
\chg Section~\ref{sec:method} (persistent acoustic graph); Section~\ref{sec:policies}; Section~\ref{sec:harmonic} (Table~\ref{tab:policyh}); Section~\ref{sec:backbones}.

\rc{R2.3}{Open science and the role of generative AI: authors must clearly delineate which specific stages of the analysis were performed by AI and ensure that all raw code and experimental logs are publicly available, allowing replication without the same AI tools.}
\resp We fully agree. The manuscript now contains a dedicated subsection with a stage-by-stage table of what the author and the AI assistant did, including design, code, running experiments, interpretation, the proof and the text; the Acknowledgments were aligned with it. The code, configurations, job lists, training and evaluation logs (including those of the diverged runs), per-recording results, trained weights and data-split manifests are public at \url{https://github.com/redwanUGA/event-aware-agn-net} (release \texttt{v1.1-revision1}). Every table, figure and quoted model result in the paper and in this letter is written into the documents from the result files by a script, apart from a few rounded summary comparisons in the text, so the tables and figures can be regenerated from the repository without any AI tool.
\chg New Section~\ref{sec:genai} (Table~\ref{tab:genai}); Data Availability Statement; Acknowledgments.

\bigskip
Thank you again for your time and for comments that have substantially strengthened the paper.

Sincerely,\\
Md. Redwan Islam
\end{document}
